\section{Simulation Methodology}
\label{sec:sim-full}

This section gives the full methodology summarized in Section~\ref{sec:sim} of the paper.

\paragraph{Why simulation} The end-to-end protocol needs local model weights, which the environment used for this paper could not download. A simulation is appropriate for the remaining questions because the quantities they depend on separate into component behaviour that we can measure without a model (detection and false-rejection rates of $\Phi$, latencies, the engine's classification, retrieval ranks) and planner behaviour that we cannot; propagating the first through the control flow of the loop, while sweeping the second, shows which unmeasured parameters decide H1--H5. We built such a seeded simulation (\texttt{benchmarks/simulation}). The simulation does not execute the agent. Its outputs follow from measured and assumed inputs. Every number in this paper belongs to one of six classes: \textsc{actual} (genuinely executed or measured on the implementation), \textsc{synthetic} (produced by the simulation), \textsc{assumed} (an input value chosen by us, not measured), \textsc{derived} (computed from existing results, e.g.\ a rate estimated from actual counts), \textsc{literature} (taken from cited work), or \textsc{proposed} (behaviour of a component that is not implemented, or an evaluation not yet done). Tables and figures state the class of their values in the caption or row label; in the text, numbers in Sections~\ref{sec:results} and~\ref{sec:safety} are \textsc{actual} unless stated, and numbers in Section~\ref{sec:sim-results} are \textsc{synthetic} unless marked otherwise. An evidence ledger lists every simulated value in the main tables with its source and class.

\paragraph{Measured inputs} Without a model we measured, on the implementation (2026-10-23, 4-core Intel Xeon at 2.8\,GHz, same software versions): the decisions of $\Phi$ and of a basic success detector (error check $E$ only) on the 92 gate states; the latency of navigation, clicks, form filling, DOM extraction, screenshots and $\Phi$ on the self-hosted sites (30 repetitions each); the failure class and strategy sequence that the real recovery engine assigns to six failure messages; and the ranks at which the real Chroma memory store returns relevant records for ten queries over a 22-record synthetic memory corpus. Detection rates used by the simulation are derived from the two held-out gate sets only (for detection, the 22 unsatisfied states left after the bot-check category is excluded, 18 from held-out~2 and 4 from held-out~1; the 24 satisfied states of the same categories for false rejection), with the estimator $(k+0.5)/(n+1)$, which gives a small nonzero rate (0.07--0.13) where none was observed, as for basic detection in five categories and $\Phi$ on extraction, and turns the 0 of 24 false rejections into 0.02. Simulated categories map to gate categories as follows: navigation, search, extraction and multi-step to themselves, data entry to text entry, and form submission to text entry and multi-step pooled. The detection rate of a missing confirmation is proxied by the multi-step rate (an assumed proxy; 3 states), and the category mapping above is likewise assumed.

\paragraph{Assumptions and generation} What needs a model---how often the planner claims completion of an unfinished task, how often actions fail and are recovered, how much memory helps---is assumed and listed with its class in Table~\ref{tab:sim-params} (Section~\ref{app:simparams}). A \emph{claim} in the simulation stands for any point at which the code evaluates $\Phi$: a \texttt{complete} action, any successful action on the last plan step (for a request that is not decomposed, every successful action), or the typing of requested exact text; $\pi$ is the probability that such a point is reached while the task is unfinished. A scenario has a task category (navigation, search, form submission, data entry, extraction, multi-step) and an outcome class (successful, incomplete, incorrect action, missing evidence, contradictory evidence, stale evidence, ambiguous) that sets the probability that the goal is actually met. Each scenario draws its random numbers once, and every configuration reuses them, so comparisons are paired by scenario. The true outcome is evaluated per configuration because memory and continuation after a rejection change the final state. For v1, parameters and code were committed (\texttt{3a0aa8e}) before the held-out set (seed 202, 100 scenarios) was generated; the reported results come from the post-hoc v2 described below; a development set (seed 101) was used only to debug the code. After the freeze, one function was changed: the latency model (Study~F) did not give all configurations the same random draws, so it was corrected to common random numbers in \texttt{41a6097}; Studies A--E and H were not affected and gave identical outputs before and after. Later, \texttt{episode()} gained three output fields (claim, true state at the claim, first decision) for the decision-level analysis; outputs were otherwise byte-identical.

\paragraph{Corrected model (v2, post hoc)} A methodology audit of the frozen model (v1) found that it contradicted the implementation in four places: its false-rejection rates came from bot-check pages, a category the scenarios do not contain; a rejected claim followed by a successful continuation was accepted without a second check, although the verification node checks every claim; the basic detector did not fail closed on an unreadable page, although it shares the error check $E$; and recovery was applied to every action failure, although action errors are first retried three times inline and a missing confirmation reaches the recovery engine only as a budget failure after 15 model calls. The results reported here come from a corrected model, v2 (\texttt{simulate\_v2.py}), that fixes these four points, draws derived rates from their Beta posteriors in the Monte Carlo replications, adds re-planning to the latency model and sweeps the continuation rate. v2 reuses v1's generator, seeds, scenario set and assumed values. It was written after the v1 results were known and is therefore \emph{not pre-registered}; the corrections moved the results for $\Phi$ in its favour (task success relative to no verification from 37 against 43, $p=0.07$, to 44 against 46; decision accuracy from 45/71 against 43/71 to 52/74 against 46/74), mainly because v1's false-rejection rate (0.089) came from bot-check pages; v1's outputs are archived (\texttt{out/heldout\_v1}) and every number that changed is listed in the statistical correction report (\texttt{artifacts/results/final\_audit/}). Monte Carlo replications regenerate the scenario set 1{,}000 times and draw the derived rates from their posteriors, to show the variability the generator and the rate uncertainty induce; the assumed values and the recovery probabilities taken from Study~E are held fixed, so the intervals omit their uncertainty.

\paragraph{Metrics} With TP, FP, FN, TN the confusion counts of reported completion against the true outcome over $N$ scenarios: $\mathrm{TSR}=\mathrm{TP}/N$; FCR and FRR as in Eq.~\eqref{eq:fcr}, so that precision $=1-\mathrm{FCR}$ and recall $=1-\mathrm{FRR}$; F1 and accuracy $(\mathrm{TP}+\mathrm{TN})/N$. Recovery success is recovered over recoverable failures; retrieval quality is precision and recall at $k=3$ and the mean reciprocal rank (MRR) of the first relevant record. Episode time is decomposed as
\begin{equation}
T = T_{\mathrm{inf}} + T_{\mathrm{browser}} + T_{\mathrm{evidence}} + T_{\mathrm{verify}} + T_{\mathrm{recovery}},
\label{eq:time}
\end{equation}
with browser, evidence and $\Phi$ times resampled from the measurements and inference time assumed. Two levels are kept apart. \emph{Policy-level} outcomes describe how an episode ends (task success, episode-level false completion); \emph{decision-level} outcomes describe whether the verifier's first decision on a completion claim matches the true state at that claim. Because a rejection can change how an episode ends, the final ground truth of a scenario may differ between configurations, so only the decision level compares verifiers against one shared truth. The decision-level outcome was defined after the v1 results were known, when an audit showed that the earlier accuracy comparison mixed the two levels. Paired configurations are compared by the exact McNemar test on per-scenario binary indicators of one level, and the variability of paired differences across regenerated scenario sets by Monte Carlo percentiles. Every $p$-value and interval computed from synthetic data is conditional on the simulation assumptions; it describes the synthetic sample, not the deployed agent.

\paragraph{Studies} Study~A generates the scenarios and computes the completion metrics, B compares four completion rules, C ablates components, D evaluates memory, E recovery, F overhead, G privacy (analytically, Section~\ref{sec:safety}) and H sensitivity.

\paragraph{Validation by recomputation} For v1, a separate AI agent of the same authoring pipeline, given only the per-scenario output files, recomputed, with separately written code, the confusion counts, rates and paired tests of Studies B and C and the memory and recovery statistics from the per-scenario files, and found no discrepancy (\texttt{simulation\_validation\_report.md}). Monte Carlo, sensitivity, latency, decision-level and recovery paired-test values were not independently recomputed. This check is a recomputation of the analysis code's arithmetic from the same synthetic files; it is neither independent human annotation of any label nor an independent replication of an experiment. For v2, a final AI review agent of the same pipeline recomputed with separately written code the policy- and decision-level confusion counts, the paired tests and the recovery and memory values from the per-scenario files, reran the simulation and the table generator byte for byte, and found no discrepancy; Monte Carlo, sensitivity and latency values were checked only through that rerun; the recomputation scripts are in \texttt{artifacts/results/final\_audit/scripts}.

